Penelitian ini bertujuan untuk mengkaji pemanfaatan pendekatan berbasis \textit{computer vision} dan \textit{deep learning} dalam mengenali pola gerakan Bahasa Isyarat Indonesia (BISINDO) secara visual. 
Fokus utama penelitian diarahkan pada pemahaman bagaimana perbedaan gestur yang tampak serupa dalam BISINDO dapat direpresentasikan dan dibedakan oleh sistem komputasi.

Secara khusus, tujuan penelitian ini adalah sebagai berikut:
\begin{enumerate}[noitemsep]
    \item Mengkaji karakteristik pola gerakan BISINDO pada tingkat kata sebagai representasi gerakan alami komunitas Tuli.
    \item Menyusun dataset berbasis \textit{landmark} yang merepresentasikan gerakan BISINDO secara terkontrol dan konsisten.
    \item Mengimplementasikan model berbasis ResNet-2D untuk mempelajari pola spasial gerakan bahasa isyarat dari data \textit{landmark}.
    \item Menilai potensi pendekatan yang diusulkan sebagai bagian dari pengembangan teknologi pendukung aksesibilitas, tanpa menggantikan peran penutur atau penerjemah manusia.
\end{enumerate}

Tujuan-tujuan tersebut dirumuskan untuk mendukung pengembangan sistem bantu yang selaras dengan praktik BISINDO dan kebutuhan komunitas Tuli.
